\uuid{DnD5}
\exo7id{540}
\titre{exo7 540}
\auteur{cousquer}
\organisation{exo7}
\datecreate{2003-10-01}
\isIndication{false}
\isCorrection{false}
\chapitre{Suite}
\sousChapitre{Suite définie par une relation de récurrence}
\cpge{mpsi-mp2i.nombres-reels-suites.i, pcsi.nombres-reels-suites.g, ptsi.nombres-reels-suites.g, mpsi-mp2i.nombres-reels-suites.d, pcsi.nombres-reels-suites.c, ptsi.nombres-reels-suites.c, mpsi-mp2i.nombres-reels-suites.e, pcsi.nombres-reels-suites.d, ptsi.nombres-reels-suites.d, bcpst1.suites-reelles, bcpst1.limites-continuite.a}
\module{Analyse}
\niveau{L1}
\difficulte{}

\contenu{
\texte{
Soit $a\in \mathbb{R}$. On considère la suite $(u_n)$ définie par 
$u_0=a$ et $u_{n+1}=e^{u_n}-2$ pour $n\geq 0$.
}
\begin{enumerate}
    \item \question{Étudier cette suite si $a=0$.}
    \item \question{Étudier cette suite si $a=-10$.}
    \item \question{Étudier cette suite si $a=3$.}
    \item \question{Généraliser en discutant selon la valeur de~$a$.}
\end{enumerate}
}
